\documentclass[10pt,a4paper]{amsart}
\usepackage[margin=19mm]{geometry}
\usepackage{amsmath,amssymb,mathtools}
\usepackage[scr=rsfs,cal=euler]{mathalfa}
\usepackage{xfrac}
\usepackage[unicode,hidelinks,bookmarks=false]{hyperref}
\newcommand{\ie}{i.e.\ }
\newcommand{\cf}{cf.\ }
\newcommand{\resp}{respectively\ }
\newcommand{\Subsection}{\subsection}
\providecommand{\nobreakdash}{\nobreak\hbox{-}}
\setlength{\emergencystretch}{2em}
\multlinegap0pt
\usepackage{siunitx}
\usepackage{siunitx}
\usepackage{siunitx}
\usepackage{siunitx}
\usepackage{amssymb}
\usepackage{mathtools}
\usepackage{siunitx}
\usepackage{xcolor}
\usepackage{tikz}
\usepackage[shortlabels]{enumitem}
\newtheorem{proposition}{Proposition}
\def\CC{{\mathbb{C}}}
\def\PP{{\mathbb{P}}}
\newcommand{\del}[1]{\frac{\partial}{\partial #1}}
\newcommand{\logg}[3]{{\rm Log}(#1;#2,#3)}
\newcommand{\mf}[1]{\mathfrak{#1}}
\newcommand{\opn}{{\mathcal O}_{\mathbb{P}^{n}}}
\newcommand{\pn}{{\mathbb{P}^{n}}}
\DeclareMathOperator{\rad}{rad}
\newcommand{\rat}[2]{{\rm Rat}(#1;#2)}
\newcommand{\sF}{\mathscr{F}}
\newcommand{\tF}{{\rm T}\sF}
\newcommand{\tpn}{T{\mathbb{P}^{n}}}
\title{Holomorphic foliations of degree two and arbitrary dimension}
\author{Maur\'icio Corr\^ea, Alan Muniz}
\date{Source: arXiv:2207.12880v5 (CC BY 4.0)}
\begin{document}
\maketitle

\noindent\emph{Source and licence.} Holomorphic foliations of degree two and arbitrary dimension. Version arXiv:2207.12880v5 (CC BY 4.0), distributed by arXiv under the Creative Commons Attribution 4.0 International licence, http://creativecommons.org/licenses/by/4.0/, as recorded in the arXiv licence metadata for this exact version and shown on the abstract page https://arxiv.org/abs/2207.12880v5. Full licence notice: \texttt{licenses/arxiv-2207.12880-CC-BY-4.0.txt}.\par\smallskip

\noindent\textit{Editorial scope.} Counted unit: the article's proposition prop:class-om-diag (source0.tex lines 980--988). The article's complete original proof of it is delivered with it (source0.tex lines 991--1111, one \texttt{proof} environment of 121 lines). No mathematics is written, repaired or re-derived: the statement and the proof are the article's own lines, in the article's own notation, and no retained line is substituted or re-flowed.  No further source-local context is needed: every cross-reference of every retained line resolves inside the two delivered spans.  Nothing was omitted from any retained span, so the package carries no omission record.  No retained line contains a citation, so the wrapper carries no bibliography and no bibliography entry.  No retained line contains a figure environment, an image-inclusion command or drawing code, so no image is delivered and the package declares no asset dependency.  Editorial formatting only, in addition: an AMS \texttt{amsart} preamble that restates the article's own declarations and macros used by the retained lines, namely the environments \texttt{proposition} on separate counters, together with the standard packages those lines need.  The retained environments print in this document as Proposition 1 in order of appearance, whereas the article prints the same items with the numbering of its own full document.  The complete native source of the article as distributed in the arXiv e-print for this exact version is bundled unchanged as \texttt{sources/128743/source0.tex}.
\begin{proposition}\label{prop:class-om-diag}
 Let $\omega\in H^0(\Omega_{\pn}^{n-2}(n+1))$ not vanishing in codimension one and let $v = x_0\del{ x_0} + \mu x_1\del{ x_1}\in H^0(\tpn)$ be such that $\iota_v \omega = 0$ and $\iota_vd\omega = a \omega$ for $a\in\CC$ a constant. Then $\omega$ defines a foliation $\sF$ of one of the following types:
 \begin{enumerate}
 \item $a = \mu + 1$, $\mu \neq -1$ and $\sF \in \logg{1^{n+1}}{2}{n}$;
 \item $a = 0$, $\mu = -1$ and $\sF\in \logg{1^{n-1},2}{2}{n}$;
 \item $a = 0$, $\mu = -2$ and $\sF \in \rat{1^{n-2},3}{n}$;
 \item $\mu \in \{1,2\}$ and $a = \mu + 2$, and $\sF$ is induced by an action of $\mf{aff}(\CC)$.
\end{enumerate}
\end{proposition}

\begin{proof}
 Since $\iota_v\omega = 0$ we may write  
 \[
 \omega = (x_0 dx_1 - \mu x_1 dx_0) \wedge \iota_{\rad}\beta + (\mu-1)x_0x_1\beta \, + \, \gamma,
 \]
 where $\beta$ and $\gamma$ do not depend on $dx_0$ nor $dx_1$ and $\iota_{\rad}\gamma = 0$. Then $\beta$ and $\gamma$ are $(n-2)$-forms depending on $n-1$ basic differentials, which means that there exists a degree-two homogeneous polynomial $q$ and a linear vector field $u = \sum_{j=2}^n b_j \del{ x_j}$ such that $\beta = \iota_u \Theta$ and $\gamma = q \, \iota_{\rad}\Theta$, where $\Theta = dx_2 \wedge \dots \wedge dx_n$. Hence,
 \[
 \omega = (x_0 dx_1 - \mu x_1 dx_0) \wedge \iota_{\rad}\iota_{u} \Theta + (\mu-1)x_0x_1\iota_u \Theta \, + \,q \, \iota_{\rad}\Theta.
 \]
 We might assume that $u$ is not proportional to the radial in the last $n-1$ variables, i.e., $u\wedge \sum_{j=2}^n x_j \del{ x_j} \neq 0$, otherwise $\omega$ would vanish in codimension one. 
 
 By hypothesis, there exists $a\in \CC$ such that $\iota_v d\omega = a \omega$. Computing the LHS, we will see that there are few possibilities for $a$. Start with expanding 
 \begin{align*}
 d\omega & = (\mu+1)dx_0 \wedge dx_1 \wedge \iota_{\rad}\iota_u \Theta - (x_0 dx_1 - \mu x_1 dx_0) \wedge d\iota_{\rad}\iota_u \Theta +\\ & \quad + (\mu-1)(x_0dx_1+x_1dx_0)\wedge \iota_u \Theta + (\mu-1)x_0x_1d\iota_u \Theta + dq\wedge \iota_{\rad}\Theta + \, q \, d\iota_{\rad}\Theta 
 \end{align*}
 then, 
 \begin{align*}
 \iota_vd\omega & = (\mu+1)(x_0 dx_1 - \mu x_1 dx_0)\wedge \iota_{\rad}\iota_u \Theta + (x_0 dx_1 - \mu x_1 dx_0) \wedge \iota_v d\iota_{\rad}\iota_u \Theta + \\ & \quad + (\mu-1)(\mu+1)x_0x_1\iota_u \Theta +(\mu-1)x_0x_1 \iota_v d\iota_u \Theta + v(q)\iota_{\rad}\Theta \\
 & = (x_0 dx_1 - \mu x_1 dx_0)\wedge ((\mu+1)\iota_{\rad}\iota_u \Theta + \iota_v d\iota_{\rad}\iota_u \Theta)+ \\ & \quad + (\mu-1)x_0x_1( (\mu+1)\iota_u \Theta + \iota_v d\iota_u \Theta) + v(q)\iota_{\rad}\Theta 
 \end{align*}
 and by $\iota_v d\omega = a \omega$ imposing we get that 
 \begin{align*}
 \iota_v d\iota_{\rad} \iota_u \Theta &= (a-\mu-1)\iota_{\rad} \iota_u \Theta, \\
 \iota_v d\iota_u \Theta &= (a-\mu-1) \iota_u \Theta + c\, \iota_{\rad} \Theta \\
 v(q) &= a \, q -(\mu-1)c \, x_0x_1
 \end{align*}
 where $c\in \CC$ and the second equation follows from the first. 
 
 Notice that $\iota_{\rad}\iota_u \Theta = \sum_{i,j} p_{ij} dx_2\wedge \dots \wedge \widehat{dx_i} \wedge \widehat{dx_j} \wedge \dots \wedge dx_n $, where each $p_{ij}$ is a quadratic polynomial not having monomials $x_0^2$, $x_0x_1$ nor $x_1^2$. Moreover,
 \[
 \iota_v d\iota_{\rad} \iota_u \Theta = \sum_{i,j} v( p_{ij} ) dx_2\wedge \dots \wedge \widehat{dx_i} \wedge \widehat{dx_j} \wedge \dots \wedge dx_n
 \]
 and we only need to compute the action of $v$ on degree two polynomials. From $v(x_0) = x_0$, $v(x_1) = \mu x_1$ and $v(x_j) = 0$ for $j\geq 2$ we get that 
 \begin{equation} \label{eq:v-action-quad}
 \frac{v(x_ix_j)}{x_ix_j} = \begin{cases} 2, & i=j = 0\\ 2\mu , & i=j=1\\
 (\mu+1) , &i = 0, \, j = 1 \\ 
 1, & i = 0, \, j\geq 2 \\ \mu ,& i = 1,\, j \geq 2 \\ 0 , & i,j\geq 2 
 \end{cases} 
 \end{equation}
 Then, to have $\iota_v d\iota_{\rad}\iota_u \Theta= (a - \mu - 1)\iota_{\rad}\iota_u \Theta$, we need $a-\mu-1\in \{ 0, 1, \mu \}$. 
\\
 \paragraph{\textbf{Case 1.0:}} Suppose that $a-\mu-1 = 0$ and that $\mu \not\in \{ -1,1\}$. Then $\iota_v d \iota_u \Theta = c \, \iota_{\rad} \Theta$ implies $c =0$. Thus $v(q) = (\mu+1)q$. From \eqref{eq:v-action-quad} and $\mu \not\in \{ -1,1\}$, $q = c' x_0x_1$ for $c'\in \CC$ a constant.
 Then 
 \[
 \omega = (x_0 dx_1 - \mu x_1 dx_0) \wedge \iota_{\rad}\iota_u \Theta + x_0x_1((\mu-1)\iota_u \Theta \, + c'\, \iota_{\rad}\Theta). 
 \]
 Consider $w = c'\, x_1 \del{ x_1} + u$ and notice that 
 \[
 \iota_w\iota_v\iota_{\rad} (dx_0\wedge \dots \wedge dx_n) = \iota_w\iota_v\iota_{\rad} (dx_0\wedge dx_1 \wedge \Theta) = \omega. 
 \]
 Also, since $u$ does not depend on $x_0$ nor $x_1$, $[v,w] =0$. Thus $\omega$ defines a foliation $\sF$ on $\pn$ with trivial tangent sheaf $\tF = \opn^{\oplus 2}$. 

 \paragraph{\textbf{Case 1.1:}} Suppose that $a-\mu-1 = 0$ and that $\mu =1$. As in the previous case, $c=0$ and $u$ does not depend on $x_0$ nor $x_1$. But now $v(q) = 2 q$ implies that $q = q_{0}x_0^2 + 2q_{1}x_0x_1+ q_{2}x_1^2$. Then
 \[
 \omega = (x_0 dx_1 - x_1 dx_0) \wedge \iota_{\rad}\iota_u \Theta + (q_{0}x_0^2 + 2q_{1}x_0x_1+ q_{2}x_1^2) \iota_{\rad}\Theta.
 \]
 Consider the vector field
 \[
 w = \left( -q_1x_0 -q_2 x_1 \right)\del{ x_0} + \left(q_0x_0 +q_1x_1 \right)\del{ x_1} + u.
 \]
 As in the previous case, we have that $ \omega = \iota_w \iota_v\iota_{\rad} (dx_0\wedge dx_1 \wedge \Theta) $, and it defines a foliation whose tangent sheaf is trivial.
\\
 \paragraph{\textbf{Case 1.2:}} Suppose that $a-\mu-1 = 0$ and that $\mu =-1$. Again, we have $c=0$ but now $v(q) = 0$, which implies that $q = q_0 x_0x_1 + q'(x_2, \dots, x_n)$. Then
 \[
 \omega = d(x_0x_1) \wedge \iota_{\rad}\iota_u \Theta -2x_0x_1\iota_u \Theta \, + \,q \, \iota_{\rad}\Theta.
 \]
 Then $\omega$ defines a foliation on $\pn$ which is the pullback of the foliation on $\PP(2,1^{(n-1)})$ given by the vector field
 \[
 m = (q_0 y_1 + q'(y_2, \dots, y_n))\del{ y_1} + \sum_{j=2}^n b_j(y)\del{ y_j}
 \]
 via the map defined by $y_1 = x_0x_1$ and $y_j = x_j$ for $j\geq 2$; recall that $u = \sum_{j=2}^n b_j(x)\del{ x_j}$.
 \\
 \paragraph{\textbf{Case 2.0:}} Suppose that $a-\mu-1 = 1$ and that $\mu \not\in \{ -2,-1,2\}$. Then $\iota_v\iota_{\rad}\iota_u \Theta = \iota_{\rad}\iota_u \Theta$ and $x_0$ must divide $\iota_{\rad}\iota_u \Theta$. Hence, there exist $\xi$ and $c'$ such that $\iota_u \Theta = x_0 \xi + c'\iota_{\rad}\Theta$. Note that $c' = -c$. Then
 \[
 \omega = x_0(x_0 dx_1 - \mu x_1 dx_0) \wedge \iota_{\rad}\xi + (\mu-1)x_0^2x_1 \xi + (-c(\mu-1)x_0x_1 + q ) \iota_{\rad}\Theta.
 \]
 On the other hand, $v(q) = (\mu+2)q -(\mu-1) c\, x_0x_1$ implies that $q = c(\mu-1)\, x_0x_1$. Hence, $\omega$ vanishes along $V(x_0)$ which contradicts our hypothesis. 
\\
 \paragraph{\textbf{Case 2.1:}} Suppose that $a-\mu-1 = 1$ and that $\mu = -2$. Then, $v(q) = 3 c x_0x_1$ which implies that $q = -3cx_0x_1 + q'(x_2, \dots, x_n)$. Hence,
 \begin{align*}
 \omega &= x_0(x_0 dx_1 +2 x_1 dx_0) \wedge \iota_{\rad}\xi -3x_0^2x_1 \xi + (3c\,x_0x_1 + q ) \iota_{\rad}\Theta = \\
 & = x_0 [(x_0 dx_1 +2 x_1 dx_0) \wedge \iota_{\rad}\xi -3x_0x_1\xi] + q' \iota_{\rad} \Theta. 
 \end{align*}
 Note that $\xi$ is a constant $(n-2)$-from on $\{dx_2, \dots, dx_n\}$. Then, up to a linear change on the last $n-1$ coordinates, we may assume that $\xi = dx_3\wedge \dots \wedge dx_n$. Then
 \[
 \omega = \iota_{\rad}( (q'dx_2- d(x_0^2x_1))\wedge dx_3 \wedge \dots \wedge dx_n).
 \]
 Now consider $f(x_2, \dots, x_n)$ a degree $3$ polynomial such that $\frac{\partial f}{\partial x_2} = q'$. It follows that 
 \[
 \omega = \iota_{\rad}( d(f-x_0^2x_1)\wedge dx_3 \wedge \dots \wedge dx_n),
 \]
 and $\omega$ defines a rational foliation of type $(3,1^{(n-2)})$. 
 \\
 \paragraph{\textbf{Case 2.2:}} Suppose that $a-\mu-1 = 1$ and that $\mu = 2$. Then, $v(q) = 4q -cx_0x_1$ which implies $q = q_0x_1^2 +c\,x_0x_1$. Then
 \begin{align*}
 \omega &= x_0(x_0 dx_1 - 2x_1 dx_0) \wedge \iota_{\rad}\xi + x_0^2x_1 \xi + (-cx_0x_1 + q ) \iota_{\rad}\Theta\\
 & = x_0(x_0 dx_1 - 2x_1 dx_0) \wedge \iota_{\rad}\xi + x_0^2x_1 \xi + q_0x_1^2 \iota_{\rad}\Theta
 \end{align*}
 As in the previous case, we may assume $\xi = dx_3\wedge \dots\wedge dx_n$, hence 
 \begin{align*}
 \omega &= x_0(x_0 dx_1 - 2x_1 dx_0) \wedge \iota_{\rad}\xi + x_0^2x_1 \xi + q_0x_1^2(x_2\xi -dx_2\wedge \iota_{\rad}\xi) \\
 & = (x_0(x_0 dx_1 - 2x_1 dx_0) - q_0x_1^2 dx_2)\wedge \iota_{\rad}\xi + (x_0^2x_1 + q_0x_1^2x_2)\xi \\
 & = \iota_{\rad}( ( x_0^2 dx_1 - 2x_1x_0 dx_0 - q_0x_1^2 dx_2) \wedge dx_3 \wedge \dots \wedge dx_n).
 \end{align*}
 Note that $q_0 \neq 0$ unless $\omega$ vanishes along $V(x_0)$. Define the vector field $w = \frac{q_0}{2}x_1 \del{ x_0} + x_0 \del{ x_2}$. It follows that 
 \[
 \omega = \iota_{\rad}\iota_w \iota_v ( ( x_0^2 dx_1 - 2x_1x_0 dx_0 - q_0x_1^2 dx_2) \wedge dx_3 \wedge \dots \wedge dx_n).
 \]
 Also notice that $[v,w] = w$ hence $\omega$ defines a foliation associated with an algebraic action of the affine Lie algebra.
\\
 \paragraph{\textbf{Case 2.3:}} Suppose that $a-\mu-1 = 1$ and that $\mu = -1$. Again $x_0$ divides $\iota_{\rad}\iota_u \Theta$. But now $v(q) = q + 2c\, x_0x_1$, which implies that $x_0$ divides $q$. Then $\omega$ vanishes along $V(x_0)$ which is absurd. 
 \\
 \paragraph{\textbf{Case 2.4:}}
 Suppose that $a-\mu-1 = 1$ and that $\mu = 1$. Then $v(q) = 3q$ implies $q=0$. Then, 
 \[
 \omega = \iota_{\rad}\iota_v\iota_u (dx_0\wedge \dots \wedge dx_n). 
 \]
 On the other hand, $\iota_vd\iota_{\rad}\iota_u\Theta = \iota_{\rad}\iota_u\Theta$ implies that $u = u' + c'\sum_{j= 2}x_j \del{ x_j}$, where the coefficients of $u'$ depend uniquely on $x_0$ and $x_1$. Note that we can replace $u$ with $u'$ in the definition of $\omega$. Moreover, $[v,u']=u'$ and $\omega$ defines a foliation tangent to an action of $\mathfrak{aff}(\CC)$.
 \\
 \paragraph{\textbf{Case 3:}} Suppose that $a-\mu-1 = \mu$. We may exchange $x_0$ and $x_1$ so that we get new parameters $(\mu',a')= (1/\mu, a/\mu)$. For $a = 2\mu + 1$ we get $a' = \mu'+2$. Then we reduce to one of the previous cases. 
\end{proof}

\end{document}
